%!TeX root=../wowtop.tex

\ArtChapter[How I Reached Home]{7head}


\lettrine[lines=4,findent=2pt]{F}{or} my own part, I remember nothing of my flight except the stress of blundering against trees and stumbling through the heather. All about me gathered the invisible terrors of the Martians; that pitiless sword of heat seemed whirling to and fro, flourishing overhead before it descended and smote me out of life. I came into the road between the crossroads and Horsell, and ran along this to the crossroads.

At last I could go no further; I was exhausted with the violence of my emotion and of my flight, and I staggered and fell by the wayside. That was near the bridge that crosses the canal by the gasworks. I fell and lay still.

I must have remained there some time.

I sat up, strangely perplexed. For a moment, perhaps, I could not clearly understand how I came there. My terror had fallen from me like a garment. My hat had gone, and my collar had burst away from its fastener. A few minutes before, there had only been three real things before me—the immensity of the night and space and nature, my own feebleness and anguish, and the near approach of death. Now it was as if something turned over, and the point of view altered abruptly. There was no sensible transition from one state of mind to the other. I was immediately the self of every day again—a decent, ordinary citizen. The silent common, the impulse of my flight, the starting flames, were as if they had been in a dream. I asked myself had these latter things indeed happened? I could not credit it.

I rose and walked unsteadily up the steep incline of the bridge. My mind was blank wonder. My muscles and nerves seemed drained of their strength. I dare say I staggered drunkenly. A head rose over the arch, and the figure of a workman carrying a basket appeared. Beside him ran a little boy. He passed me, wishing me good night. I was minded to speak to him, but did not. I answered his greeting with a meaningless mumble and went on over the bridge.

Over the Maybury arch a train, a billowing tumult of white, firelit smoke, and a long caterpillar of lighted windows, went flying south—clatter, clatter, clap, rap, and it had gone. A dim group of people talked in the gate of one of the houses in the pretty little row of gables that was called Oriental Terrace. It was all so real and so familiar. And that behind me! It was frantic, fantastic! Such things, I told myself, could not be.

Perhaps I am a man of exceptional moods. I do not know how far my experience is common. At times I suffer from the strangest sense of detachment from myself and the world about me; I seem to watch it all from the outside, from somewhere inconceivably remote, out of time, out of space, out of the stress and tragedy of it all. This feeling was very strong upon me that night. Here was another side to my dream.

But the trouble was the blank incongruity of this serenity and the swift death flying yonder, not two miles away. There was a noise of business from the gasworks, and the electric lamps were all alight. I stopped at the group of people.

»What news from the common?« said I.

There were two men and a woman at the gate.

»Eh?« said one of the men, turning.

»What news from the common?« I said.

»Ain't yer just \textit{been} there?« asked the men.

»People seem fair silly about the common,« said the woman over the gate. »What's it all abart?«

»Haven't you heard of the men from Mars?« said I; »the creatures from Mars?«

»Quite enough,« said the woman over the gate. »Thenks«; and all three of them laughed.

I felt foolish and angry. I tried and found I could not tell them what I had seen. They laughed again at my broken sentences.

»You'll hear more yet,« I said, and went on to my home.

I startled my wife at the doorway, so haggard was I. I went into the dining room, sat down, drank some wine, and so soon as I could collect myself sufficiently I told her the things I had seen. The dinner, which was a cold one, had already been served, and remained neglected on the table while I told my story.

»There is one thing,« I said, to allay the fears I had aroused; »they are the most sluggish things I ever saw crawl. They may keep the pit and kill people who come near them, but they cannot get out of it\textellipsis But the horror of them!«

»Don't, dear!« said my wife, knitting her brows and putting her hand on mine.

»Poor Ogilvy!« I said. »To think he may be lying dead there!«

My wife at least did not find my experience incredible. When I saw how deadly white her face was, I ceased abruptly.

»They may come here,« she said again and again.

I pressed her to take wine, and tried to reassure her.

»They can scarcely move,« I said.

I began to comfort her and myself by repeating all that Ogilvy had told me of the impossibility of the Martians establishing themselves on the earth. In particular I laid stress on the gravitational difficulty. On the surface of the earth the force of gravity is three times what it is on the surface of Mars. A Martian, therefore, would weigh three times more than on Mars, albeit his muscular strength would be the same. His own body would be a cope of lead to him, therefore. That, indeed, was the general opinion. Both \textit{The Times} and the \textit{Daily Telegraph}, for instance, insisted on it the next morning, and both overlooked, just as I did, two obvious modifying influences.

The atmosphere of the earth, we now know, contains far more oxygen or far less argon (whichever way one likes to put it) than does Mars'. The invigorating influences of this excess of oxygen upon the Martians indisputably did much to counterbalance the increased weight of their bodies. And, in the second place, we all overlooked the fact that such mechanical intelligence as the Martian possessed was quite able to dispense with muscular exertion at a pinch.

But I did not consider these points at the time, and so my reasoning was dead against the chances of the invaders. With wine and food, the confidence of my own table, and the necessity of reassuring my wife, I grew by insensible degrees courageous and secure.

»They have done a foolish thing,« said I, fingering my wineglass. »They are dangerous because, no doubt, they are mad with terror. Perhaps they expected to find no living things—certainly no intelligent living things.«

»A shell in the pit,« said I, »if the worst comes to the worst, will kill them all.«

\begin{wrapfigure}{O}{0.5\textwidth}
\centering
\includegraphics[width=0.5\textwidth]{7tailpiece}
\end{wrapfigure}

The intense excitement of the events had no doubt left my perceptive powers in a state of erethism. I remember that dinner table with extraordinary vividness even now. My dear wife's sweet anxious face peering at me from under the pink lamp shade, the white cloth with its silver and glass table furniture—for in those days even philosophical writers had many little luxuries—the crimson-purple wine in my glass, are photographically distinct. At the end of it I sat, tempering nuts with a cigarette, regretting Ogilvy's rashness, and denouncing the short-sighted timidity of the Martians.

So some respectable dodo in the Mauritius might have lorded it in his nest, and discussed the arrival of that shipful of pitiless sailors in want of animal food. »We will peck them to death tomorrow, my dear.«

I did not know it, but that was the last civilised dinner I was to eat for very many strange and terrible days.

%\begin{figure}[b!]
%\centering
%\includegraphics[width=\tailpiecesize]{7tailpiece}\captionlistentry{Tailpiece to Chapter \thechapter}
%\end{figure}